\chapter{84}

\section{Bettmeralp/VS (CH), Freitag, 7.
September}



Manuela und Ava übernahmen die Küchenarbeit. Joh erledigte Bürokram. Und
Louis entschuldigte sich mit dringend anstehenden Anrufen.

\sectionbreak

In überschaubarer Nähe oberhalb des Hauses, wenige Schritte weiter Richtung Norden, entdeckte Louis eine Bank.
Er wischte die abgeblätterte Farbe weg, setzte sich und streckte Arme und Beine von sich.
Er seufzte. \emph{Che Casino}\footnote{Was für ein Theater.}. Wie würde diese Katastrophe nur enden?
Es wurde nur schlimmer.
Louis legte sich die Hände übers Gesicht und stöhnte leise vor sich hin. 
Wenn Bastiano bei Valentina seine Nummer gesehen hatte, dann war er
geliefert. Es war wahrscheinlich. Ebenso, dass er ihr Gespräch
mit angehört hatte. Wie kam der sonst auf den Gedanken, Louis sei in der
Nacht bei Giorgia am Felsen gewesen? Oder bluffte er und wollte
Valentina unter Druck setzen? Wollte er sie für sich gewinnen? Sie gegen
ihn, Louis, aufhetzen? Er musste Valentina noch einmal anrufen. Er hatte
sie vorhin einfach weggedrückt.
Louis stand auf und ging vor der Bank hin und her.
Vielleicht sollte er endgültig jemanden um Hilfe bitten. Marco? Doch bei
ihm hatte er sich nicht mehr gemeldet. Oder Maman? Oder Lucia? Jetzt
plötzlich? Und was erwartete er von ihnen?
Er zog das Handy aus der Hosentasche. Als Erstes wollte er nach dem Video suchen,  
das Valentina erwähnt hatte. Der vermummte Mann, der in der
Nacht, die \emph{Questur} gestürmt hatte, sei darauf zu sehen. Der
wurde gesucht. Warum, war Louis noch immer schleierhaft. Der Aufruf
musste im Internet zu finden sein. Louis scrollte durch die
verschiedenen Nachrichtenportale. Vielleicht gab es einen Hinweis in
einer der grossen Tageszeitungen. Er stutzte. Das musste es sein. In
der \emph{«Corriere della Sera»}, 6. September, zweite Seite, gross
aufgemacht, fand er ein Bild. Es war verschwommen. Die Augenpartie verpixelt.
Darunter eine Mitteilung der Mailänder Polizei.

\sectionbreak

\indentedblock{«Nächtlicher Auftritt des Mannes auf der Polizeistation geklärt

Der Mann, der in den frühen Morgenstunden des 20. August vermummt auf
der Polizeistation Milano aufgetaucht ist, konnte identifiziert werden.
Ob er mit dem Sturz der jungen Mailänderin vom 20. August am «Roccia del
Culto» im Zusammenhang steht, ist noch nicht geklärt. Aus
ermittlungstechnischen Gründen gelangen aktuell keine weiteren Details
an die Öffentlichkeit. Die Polizei bedankt sich bei der Bevölkerung für
ihre Mithilfe. Il Commissario Alessandro Mantovani, Polizia di Stato
Milano.»}

\sectionbreak

Ob Valentina die Meldung gesehen hatte?
Louis’ Handy vibrierte. Valentina? Er hielt den Atem an. Ein anonymer Anruf?
Er starrte auf das Display. Das Klingeln verstummte.
Louis drückte das Gerät mit aller Kraft zusammen, streckte es in die
Luft und biss die Zähne zusammen. \emph{Merda, merda}.\footnote{Scheisse, Scheisse.}
Ein zertrümmertes Handy löste nichts. Gar nichts. Er durfte sich nicht
verrückt machen. Anonyme Anrufe gab es immer mal wieder. Das musste
nichts heissen. Und wenn - er senkte den Arm. Ein Schluchzen begann ihn
zu schütteln. Wieder erschrak er über seine abgründige Verzweiflung, oder eigentlich viel mehr 
über sein Weinen, das er nicht steuern und noch weniger unterdrücken konnte. 
Wäre das nicht genug, folgte seinen Tränen postwendend ein Gefühl der Beschämung. 
Ein selbstmitleidiges Kind, ein Jammerlappen war er.  
Er brauchte dringend weiterhin eine Verbindung zu Mailand.
Wenigstens eine zu Valentina. Louis wischte sich die feuchten Wangen trocken. 
Eine bleierne Müdigkeit hielt ihn auf die Bank gedrückt.
Seine Augen fielen immer wieder zu. Er brauchte Schlaf. Weitere Entscheidungen mussten bis Morgen warten.

Louis betrat leise die \emph{Schäferloft}. Ava schlief.
